\chapter*{序言}

本书是北京大学通用人工智能实验班（简称“通班”）的大二专业必修课程《AI中的数学》的教材. 随着时代的发展，相比传统的机器学习，AI技术已经有了颠覆性的变化. 深度学习、强化学习、生成式大模型等新兴技术的出现，使得AI在各个领域的应用变得更加广泛和深入. 然而，经典的AI或者机器学习教材往往侧重于传统方法，缺乏对现代AI技术所需数学基础的系统介绍，更缺乏对于这些数学背后思想的深入探讨. 因此，我们有了编写本书的想法. 

遗憾的是，现代的AI社区，理论与实践的讨论日益割裂，深度学习理论大都远远落后于实践的发展. 这些理论不但越来越复杂，而且往往依赖过于理想化、甚至不可验证的假设，使得理论结果难以应用于实际问题. 因此，我们认为，AI中真正有用的数学绝对不是来自这些复杂的理论. 相反，对现代AI发展做出突出贡献的这些工作，都是来自那些受到良好数学训练的研究者们，利用这些数学的直觉和思想所发展出来的. 本书的一大目标，就是努力挖掘这些工作背后的数学直觉和思想，然后把它们整理成系统的内容，帮助读者建立起有用的数学直觉. 

“AI中的数学”这门课，从2022年开始在北京大学通班开设，到今天已经有四届学生学习过这门课. 很明显，这门课程不同于绝大部分传统的数学课程，也不同于绝大部分的AI课程，同学们在学习的过程中既有新奇的体验，也有全新的挑战. 这门课内容多且杂，充满思想和哲学的讨论，而且对于大二的学生来说，缺乏真正的AI实践，这些全新的概念和思想也并不容易理解，往往就是囫囵吞枣地过了一遍，难以消化吸收. 也有很多同学非常喜欢这门课的内容，同样是因为它这样地与众不同. 因此，这门课，或者说这本书，成为了一个饱受争议的对象. 

然而，经过几年的回访和调研，这本书中的内容，正如我们的目标，在这些同学后续的学习和研究中，发挥了重要的作用. 无论是在学术界进行AI相关的研究，还是在工业界从事AI相关的工作，这些同学都表示，每当遇到无法理解和解决的数学问题时，他们真的会求助这本书，然后理解当年没学懂的东西. 因此，我们感到十分幸运，这本书对不同层次的读者都可以有所收获. 同样的理由，即便这本书对于初学者饱受争议，我们依然决定以这样的方式呈现这些内容：不求一次学明白，而是把这些思想留给读者，随着经验的积累不断深化自己的认知. 

本书的选材包罗万象，但都以一条非常清晰的脉络串联而成. 书中的内容包含了过去在AI社区发挥重要作用的思想，也包括了未来我们认为会变得十分重要的思想. 十分有趣的是，本书的编撰和完善的过程，恰好见证了大模型的兴起和惊人进步的过程. 我们在这个过程中也在不断检验本书内容的有效性和实用性. 事实证明，这些数学思想和直觉，曾经可能是不重要、被认为简单的东西，在大语言模型的时代，它们变得无比重要，成为理解和改进这些模型的关键. 

这门课程也见证着AI对于教育的巨大变革. AI技术的发展，使得教育资源的获取变得更加便捷和高效，也使得教学的形式、老师的职责、考核制度都发生着巨大的变化. 在本课程中，我们大量探索了AI技术在教学中的应用，例如如何利用AI组织教学内容、辅助学生问答、批改作业，以及在这样的背景下如何设计考核. 因此，这本书的编排也充分也是考虑了AI时代自学的需要，不是传统的定理证明、定义辨析、例题等形式，因为这些东西当代的大模型都掌握得非常好；相反，我们将材料的重点放在了数学内容的动机、如何想出来这样的数学以及这样的数学和AI的关系是什么这些更加思辨、不在纸面上被着重强调的内容. 希望这样的考虑可以帮助读者更方便地进行自学. 

本书编纂的过程中，得到了许多人的帮助和支持. 在此，我们要特别感谢北京大学通用人工智能实验班的2019级到2023级的各位同学们以及来自信息科学技术学院、数学科学学院、生命科学学院、化学与分子工程学院、临床医学院、信息管理系、公共卫生学院、外语学院、光华管理学院、元培学院热情的选课同学们（写到这里才发现，难以想象有如此多的院系），他们在课程学习过程中提出了许多宝贵的意见和建议，帮助我们不断改进和完善教材内容. 同时，我们也要感谢陈昱蓉、段志健、肖茜之、孙皓然、牟湛存、潘凝心、王颖，他们是历届的助教团，对于这个课程内容的建设起到了至关重要的作用. 此外，我们还非常感谢李东晨、唐静吾、李炳辉、王彦晶，他们是本书中涉及领域的专家，对于本书内容的建议、指正和反馈，都极大地提升了本书的质量. 最后，我们要感谢北京大学人工智能研究院给我们这个机会开设这门课程，并且给了我们极大的支持和帮助. 

{
\raggedleft\textit{邓小铁\quad 李翰禹}

\raggedleft\textit{2026年1月4日\quad 于燕园}

}